% Options for packages loaded elsewhere
% Options for packages loaded elsewhere
\PassOptionsToPackage{unicode}{hyperref}
\PassOptionsToPackage{hyphens}{url}
\PassOptionsToPackage{dvipsnames,svgnames,x11names}{xcolor}
%
\documentclass[
  letterpaper,
  DIV=11,
  numbers=noendperiod]{scrartcl}
\usepackage{xcolor}
\usepackage{amsmath,amssymb}
\setcounter{secnumdepth}{-\maxdimen} % remove section numbering
\usepackage{iftex}
\ifPDFTeX
  \usepackage[T1]{fontenc}
  \usepackage[utf8]{inputenc}
  \usepackage{textcomp} % provide euro and other symbols
\else % if luatex or xetex
  \usepackage{unicode-math} % this also loads fontspec
  \defaultfontfeatures{Scale=MatchLowercase}
  \defaultfontfeatures[\rmfamily]{Ligatures=TeX,Scale=1}
\fi
\usepackage{lmodern}
\ifPDFTeX\else
  % xetex/luatex font selection
\fi
% Use upquote if available, for straight quotes in verbatim environments
\IfFileExists{upquote.sty}{\usepackage{upquote}}{}
\IfFileExists{microtype.sty}{% use microtype if available
  \usepackage[]{microtype}
  \UseMicrotypeSet[protrusion]{basicmath} % disable protrusion for tt fonts
}{}
\makeatletter
\@ifundefined{KOMAClassName}{% if non-KOMA class
  \IfFileExists{parskip.sty}{%
    \usepackage{parskip}
  }{% else
    \setlength{\parindent}{0pt}
    \setlength{\parskip}{6pt plus 2pt minus 1pt}}
}{% if KOMA class
  \KOMAoptions{parskip=half}}
\makeatother
% Make \paragraph and \subparagraph free-standing
\makeatletter
\ifx\paragraph\undefined\else
  \let\oldparagraph\paragraph
  \renewcommand{\paragraph}{
    \@ifstar
      \xxxParagraphStar
      \xxxParagraphNoStar
  }
  \newcommand{\xxxParagraphStar}[1]{\oldparagraph*{#1}\mbox{}}
  \newcommand{\xxxParagraphNoStar}[1]{\oldparagraph{#1}\mbox{}}
\fi
\ifx\subparagraph\undefined\else
  \let\oldsubparagraph\subparagraph
  \renewcommand{\subparagraph}{
    \@ifstar
      \xxxSubParagraphStar
      \xxxSubParagraphNoStar
  }
  \newcommand{\xxxSubParagraphStar}[1]{\oldsubparagraph*{#1}\mbox{}}
  \newcommand{\xxxSubParagraphNoStar}[1]{\oldsubparagraph{#1}\mbox{}}
\fi
\makeatother


\usepackage{longtable,booktabs,array}
\newcounter{none} % for unnumbered tables
\usepackage{calc} % for calculating minipage widths
% Correct order of tables after \paragraph or \subparagraph
\usepackage{etoolbox}
\makeatletter
\patchcmd\longtable{\par}{\if@noskipsec\mbox{}\fi\par}{}{}
\makeatother
% Allow footnotes in longtable head/foot
\IfFileExists{footnotehyper.sty}{\usepackage{footnotehyper}}{\usepackage{footnote}}
\makesavenoteenv{longtable}
\usepackage{graphicx}
\makeatletter
\newsavebox\pandoc@box
\newcommand*\pandocbounded[1]{% scales image to fit in text height/width
  \sbox\pandoc@box{#1}%
  \Gscale@div\@tempa{\textheight}{\dimexpr\ht\pandoc@box+\dp\pandoc@box\relax}%
  \Gscale@div\@tempb{\linewidth}{\wd\pandoc@box}%
  \ifdim\@tempb\p@<\@tempa\p@\let\@tempa\@tempb\fi% select the smaller of both
  \ifdim\@tempa\p@<\p@\scalebox{\@tempa}{\usebox\pandoc@box}%
  \else\usebox{\pandoc@box}%
  \fi%
}
% Set default figure placement to htbp
\def\fps@figure{htbp}
\makeatother





\setlength{\emergencystretch}{3em} % prevent overfull lines

\providecommand{\tightlist}{%
  \setlength{\itemsep}{0pt}\setlength{\parskip}{0pt}}



 
\usepackage[]{biblatex}


\KOMAoption{captions}{tableheading}
\makeatletter
\@ifpackageloaded{caption}{}{\usepackage{caption}}
\AtBeginDocument{%
\ifdefined\contentsname
  \renewcommand*\contentsname{Table of contents}
\else
  \newcommand\contentsname{Table of contents}
\fi
\ifdefined\listfigurename
  \renewcommand*\listfigurename{List of Figures}
\else
  \newcommand\listfigurename{List of Figures}
\fi
\ifdefined\listtablename
  \renewcommand*\listtablename{List of Tables}
\else
  \newcommand\listtablename{List of Tables}
\fi
\ifdefined\figurename
  \renewcommand*\figurename{Figure}
\else
  \newcommand\figurename{Figure}
\fi
\ifdefined\tablename
  \renewcommand*\tablename{Table}
\else
  \newcommand\tablename{Table}
\fi
}
\@ifpackageloaded{float}{}{\usepackage{float}}
\floatstyle{ruled}
\@ifundefined{c@chapter}{\newfloat{codelisting}{h}{lop}}{\newfloat{codelisting}{h}{lop}[chapter]}
\floatname{codelisting}{Listing}
\newcommand*\listoflistings{\listof{codelisting}{List of Listings}}
\makeatother
\makeatletter
\makeatother
\makeatletter
\@ifpackageloaded{caption}{}{\usepackage{caption}}
\@ifpackageloaded{subcaption}{}{\usepackage{subcaption}}
\makeatother
\usepackage{bookmark}
\IfFileExists{xurl.sty}{\usepackage{xurl}}{} % add URL line breaks if available
\urlstyle{same}
% fallback for those not using the hyperref driver hyperxmp:
\makeatletter
\@ifundefined{xmpquote}{\newcommand{\xmpquote}[1]{#1}}{}
\makeatother
\hypersetup{
  pdftitle={Lendingclub's loan default prediction},
  pdfauthor={Rafiq Islam},
  colorlinks=true,
  linkcolor={blue},
  filecolor={Maroon},
  citecolor={Blue},
  urlcolor={Blue},
  pdfcreator={LaTeX via pandoc}}


\title{Lendingclub's loan default prediction}
\author{Rafiq Islam}
\date{2023-02-23}
\begin{document}
\maketitle

\renewcommand*\contentsname{Table of contents}
{
\setcounter{tocdepth}{4}
\tableofcontents
}

\href{../../../codepages/lendingclub/index.qmd}{Notebook}

\subsection{Project Overview}\label{project-overview}

LendingClub is a U.S.-based financial services company that initially
began as a peer-to-peer (P2P) lending platform, allowing individual
investors to lend directly to individual borrowers through a
marketplace. Founded in 2006, it grew quickly to become one of the
largest and most popular P2P lending platforms, helping connect
borrowers in need of personal loans with investors looking for
alternative investment opportunities. LendingClub has been known for its
transparent data-sharing practices, making anonymized loan data
available to researchers, data scientists, and investors. This data is
widely used in financial research, especially for predictive modeling of
loan risk and borrower behavior. The aim of this data science project is
to build a machine learning model to predict the likelihood of a loan
default.

\subsubsection{More information about
LendingClub}\label{more-information-about-lendingclub}

\begin{enumerate}
\def\labelenumi{\arabic{enumi}.}
\tightlist
\item
  \textbf{P2P Lending Model (Early Focus):}

  \begin{itemize}
  \tightlist
  \item
    Borrowers could apply for loans, typically unsecured personal loans,
    through LendingClub's platform.
  \item
    Loans were then funded by individual investors who could review
    borrowers' profiles, risk grades, and other financial information
    before committing funds.
  \item
    The model offered borrowers a way to access loans outside
    traditional banks, often at lower interest rates, and allowed
    investors to diversify by spreading investments across multiple
    loans.
  \end{itemize}
\item
  \textbf{Risk and Return:}

  \begin{itemize}
  \tightlist
  \item
    LendingClub assigned credit grades (A--G) to each loan based on
    creditworthiness, which affected interest rates. Higher risk meant
    potentially higher returns for investors but also higher chances of
    default.
  \item
    Investors bore the risk if a borrower defaulted, which was a notable
    risk factor compared to FDIC-insured deposits.
  \end{itemize}
\item
  \textbf{Shift to a Bank Model:}

  \begin{itemize}
  \tightlist
  \item
    Over time, LendingClub transitioned away from P2P lending and
    restructured as a more traditional bank, obtaining a bank charter in
    2021.
  \item
    It now offers banking products, such as high-yield savings accounts,
    and operates more like a digital bank while still focusing on
    lending products.
  \end{itemize}
\item
  \textbf{Borrower and Loan Profiles:}

  \begin{itemize}
  \tightlist
  \item
    LendingClub primarily focuses on personal loans for debt
    consolidation, credit card refinancing, home improvement, and other
    purposes.
  \item
    Borrowers' profiles typically include information on income, credit
    score, debt-to-income ratio, and loan purpose, which is used for
    assessing risk.
  \end{itemize}
\end{enumerate}

\subsection{Dataset}\label{dataset}

The dataset is a publicly available data from
\href{https://www.kaggle.com/datasets/jeandedieunyandwi/lending-club-dataset}{kaggle.com}.
It originally contains 396030 entries, with 100.4 MB. However, for easy
github push, I reduce the dataset slightly in order to have size smaller
than 100 MB. So, in the reduced form, it has 395900 entries with the
following columns:

\begin{itemize}
\tightlist
\item
  \texttt{loan\_amnt:} The listed amount of the loan applied for by the
  borrower. If at some point in time, the credit department reduces the
  loan amount, then it will be reflected in this value.\\
\item
  \texttt{term:} The number of payments on the loan. Values are in
  months and can be either 36 or 60.\\
\item
  \texttt{int\_rate:} Interest Rate on the loan installment The monthly
  payment owed by the borrower if the loan originates.\\
\item
  \texttt{grade:} LC assigned loan grade\\
\item
  \texttt{sub\_grade:} LC assigned loan subgrade\\
\item
  \texttt{emp\_title:} The job title supplied by the Borrower when
  applying for the loan.\\
\item
  \texttt{emp\_length:} Employment length in years. Possible values are
  between 0 and 10 where 0 means less than one year and 10 means ten or
  more years.\\
\item
  \texttt{home\_ownership:} The home ownership status provided by the
  borrower during registration or obtained from the credit report. Our
  values are: RENT, OWN, MORTGAGE, OTHER\\
\item
  \texttt{annual\_inc:} The self-reported annual income provided by the
  borrower during registration.\\
\item
  \texttt{verification\_status:} Indicates if income was verified by LC,
  not verified, or if the income source was verified\\
\item
  \texttt{issue\_d:} The month which the loan was funded\\
\item
  \texttt{loan\_status:} Current status of the loan\\
\item
  \texttt{purpose:} A category provided by the borrower for the loan
  request.\\
\item
  \texttt{title:} The loan title provided by the borrower\\
\item
  \texttt{zip\_code:} The first 3 numbers of the zip code provided by
  the borrower in the loan application.\\
\item
  \texttt{addr\_state:} The state provided by the borrower in the loan
  application.\\
\item
  \texttt{dti:} A ratio calculated using the borrower's total monthly
  debt payments on the total debt obligations, excluding mortgage and
  the requested LC loan, divided by the borrower's self-reported monthly
  income.\\
\item
  \texttt{earliest\_cr\_line:} The month the borrower's earliest
  reported credit line was opened.\\
\item
  \texttt{open\_acc:} The number of open credit lines in the borrower's
  credit file.\\
\item
  \texttt{pub\_rec:} Number of derogatory public records.\\
\item
  \texttt{revol\_bal:} Total credit revolving balance.\\
\item
  \texttt{revol\_util:} Revolving line utilization rate, or the amount
  of credit the borrower is using relative to all available revolving
  credit.\\
\item
  \texttt{total\_acc:} The total number of credit lines currently in the
  borrower's credit file
\item
  \texttt{initial\_list\_status:} The initial listing status of the
  loan. Possible values are -- W, F
\item
  \texttt{application\_type:} Indicates whether the loan is an
  individual application or a joint application with two co-borrowers.\\
\item
  \texttt{mort\_acc:} Number of mortgage accounts.\\
\item
  \texttt{pub\_rec\_bankruptcies:} Number of public record bankruptcies
\end{itemize}

Here \texttt{loan\_status} is the predictive or dependent variable and
the rest are predicting or independent variables aka features. We want
to predict if the loan will be \texttt{Fully\ Paid} or
\texttt{Charged\ Off} given the feature values.

\subsection{Stakeholders}\label{stakeholders}

\subsection{Key Performance Indicators (KPIs) of
LendingClub}\label{key-performance-indicators-kpis-of-lendingclub}

\subsection{Modeling}\label{modeling}

\subsection{Results and Outcome}\label{results-and-outcome}

\subsection{Future Directions}\label{future-directions}


\printbibliography



\end{document}
